\section{Incidence data and reductions of the witnesses}\label{app:certificates}

The tables below specify the four graphs used in the constructions
and the finite calculations in their proofs. In particular, their
definitions do not depend on recovering edges from a drawing.

\begin{table}[htbp]
\centering\small
\caption{For each vertex $i$, the entries are precisely its neighbours
larger than $i$. Thus each edge occurs once. The symbol $\varnothing$
means that there is no larger neighbour; the dash indicates that the
vertex is absent.}\label{tab:witness-adjacency}
\begin{tabular}{@{}rllll@{}}
\toprule
$i$ & $\Gt$ & $\Ga$ & $\Gb$ & $\Gc$ \\
\midrule
0 & $4,5,9,10,11,12$ & $4,5,8,10,12,13$ & $5,7,9,12,13$ & $4,5,7,9,12,13$ \\
1 & $4,6,8,9,11,12$ & $5,7,8,10,12,13$ & $5,8,9,11,13$ & $4,5,8,10,12,13$ \\
2 & $5,7,8,9,10,12$ & $6,7,9,12,13$ & $6,7,8,11,12,13$ & $6,7,8,10,11,13$ \\
3 & $6,7,8,10,11,12$ & $6,9,10,12,13$ & $6,7,9,10,12,13$ & $6,7,9,10,11,13$ \\
4 & $8,10,12$ & $8,9,11,13$ & $6,8,9,10,11,13$ & $8,9,11,13$ \\
5 & $8,11,12$ & $9,10,11,12$ & $7,8,10,11,12$ & $9,10,11,13$ \\
6 & $9,10,12$ & $7,10,11,12$ & $10,11,12$ & $8,9,12,13$ \\
7 & $9,11,12$ & $8,11,13$ & $10,11,13$ & $10,12,13$ \\
8 & $10,11$ & $11,12$ & $10,12,13$ & $11,12$ \\
9 & $10,11$ & $11,12,13$ & $11,12,13$ & $11,12$ \\
10 & $\varnothing$ & $11,13$ & $\varnothing$ & $11,12$ \\
11 & $\varnothing$ & $\varnothing$ & $\varnothing$ & $\varnothing$ \\
12 & $\varnothing$ & $\varnothing$ & $\varnothing$ & $\varnothing$ \\
13 & --- & $\varnothing$ & $\varnothing$ & $\varnothing$ \\
\bottomrule
\end{tabular}

\end{table}

\subsection{A triangulated disc for the thirteen-vertex graph}
\label{app:disc-certificates}

Here is an incidence description of the disc in Figure~\ref{fig:G13}.
Its vertices are $0,1,\ldots,7,12$ and $b_0,\ldots,b_{11}$; its
triangles are listed in Table~\ref{tab:G13-disc-triangles}.
In the figure, $b_0$ is the topmost boundary vertex, and the boundary
indices increase clockwise.
The interior vertices keep their labels. On the boundary, the map
to $\Gt$ is
\[
 b_{4r}\longmapsto10,\qquad
 b_{4r+1}\longmapsto9,\qquad
 b_{4r+2}\longmapsto11,\qquad
 b_{4r+3}\longmapsto8\qquad(0\le r\le2).
\]

\begin{table}[htbp]
\centering\small
\caption{The triangles of the disc whose boundary is attached with
degree three.}\label{tab:G13-disc-triangles}
\begin{tabular}{@{}rlrlrlrl@{}}
\toprule
$j$ & $T_j$ & $j$ & $T_j$ & $j$ & $T_j$ & $j$ & $T_j$ \\
\midrule
1 & $\{0,4,b_{4}\}$ & 8 & $\{1,4,12\}$ & 15 & $\{2,7,b_{9}\}$ & 22 & $\{3,7,12\}$ \\
2 & $\{0,4,12\}$ & 9 & $\{1,6,b_{1}\}$ & 16 & $\{2,7,12\}$ & 23 & $\{3,b_{11},b_{0}\}$ \\
3 & $\{0,5,b_{6}\}$ & 10 & $\{1,6,12\}$ & 17 & $\{2,b_{7},b_{8}\}$ & 24 & $\{3,b_{11},b_{10}\}$ \\
4 & $\{0,5,12\}$ & 11 & $\{1,b_{3},b_{2}\}$ & 18 & $\{2,b_{9},b_{8}\}$ & 25 & $\{4,b_{3},b_{4}\}$ \\
5 & $\{0,b_{5},b_{4}\}$ & 12 & $\{1,b_{1},b_{2}\}$ & 19 & $\{3,6,b_{0}\}$ & 26 & $\{5,b_{7},b_{6}\}$ \\
6 & $\{0,b_{5},b_{6}\}$ & 13 & $\{2,5,b_{7}\}$ & 20 & $\{3,6,12\}$ & 27 & $\{6,b_{1},b_{0}\}$ \\
7 & $\{1,4,b_{3}\}$ & 14 & $\{2,5,12\}$ & 21 & $\{3,7,b_{10}\}$ & 28 & $\{7,b_{9},b_{10}\}$ \\
\bottomrule
\end{tabular}

\end{table}

To verify the disc assertion, start with $T_1$ and attach the remaining
triangles in this order:
\[
\begin{gathered}
2,4,3,5,6,8,7,10,9,11,12,14,13,16,15,17,18,\\
20,19,22,21,23,24,25,26,27,28.
\end{gathered}
\]
At each step, the intersection with the preceding union is either
one boundary edge or two consecutive boundary edges, with no
additional vertices. This follows directly from the triangle list.
Attaching a triangle along either such boundary arc preserves a
closed disc. The final boundary is
$b_0,b_1,\ldots,b_{11},b_0$.

The triangles map bijectively onto the $28$ triangles of $\cl(\Gt)$.
Their edges give exactly the adjacency data for $\Gt$, which contain
no $K_4$, so the image is the entire clique complex. Each interior
vertex has a distinct label outside $\{8,9,10,11\}$. The images of
interior edges are distinct, and each has an endpoint outside the
four-cycle. Thus only equal-labelled boundary vertices and their
corresponding boundary edges are identified. This verifies the
hypotheses of Proposition~\ref{prop:disc}, with boundary degree three.

\subsection{The obstruction to a first deletion}

The graphs in Table~\ref{tab:witness-links} are $K_4$-free and have
connected vertex neighbourhoods. The table gives their first Betti
numbers $|E(\Nb_G(v))|-|V(\Nb_G(v))|+1$ in vertex-label order.
All are positive; every edge neighbourhood is independent with at
least two vertices. Hence all these neighbourhoods have nonzero
reduced homology and none belongs to $\IS$.

\begin{table}[htbp]
\centering\small
\caption{Nonzero homology of the vertex and edge neighbourhoods.
The Betti numbers of vertex neighbourhoods have the same values over
every field.}\label{tab:witness-links}
\begin{tabular}{@{}lll@{}}
\toprule
Graph & First Betti numbers of vertex neighbourhoods & Edge-neighbourhood sizes\\
\midrule
$\Gt$ & $1,1,1,1,1,1,1,1,2,2,2,2,1$ & $2,3$\\
$\Ga$ & $1,1,1,1,1,2,1,1,1,2,2,1,2,2$ & $2,3$\\
$\Gb$ & $1,1,1,1,1,2,1,2,2,2,1,1,1,3$ & $2,3,4$\\
$\Gc$ & $2,1,1,1,1,1,1,2,1,2,2,2,1,2$ & $2,3$\\
\bottomrule
\end{tabular}
\end{table}

\subsection{Reductions of the one-edge extensions}

Tables~\ref{tab:Ga-reduction} and~\ref{tab:variant-reductions} give
the local neighbourhood at every deletion. A displayed path specifies
exactly its consecutive edges. A single label denotes a one-vertex
graph. At step $13$ for $\Gb+\{0,3\}$, the two displayed paths are
united at their common vertices; their union is a tree.
All neighbourhoods are evaluated before the corresponding deletion,
and all three reductions end at the vertex $13$.

\begin{table}[htbp]
\centering\small
\caption{The complete reduction of $\Ga+\{0,1\}$ used in
Theorem~\ref{thm:B}.}\label{tab:Ga-reduction}
\begin{tabular}{@{}rll@{}}
\toprule
Step & Delete & Current neighbourhood \\
\midrule
1 & $e(0,5)$ & $10\text{--}1\text{--}12$ \\
2 & $e(0,10)$ & $1\text{--}13$ \\
3 & $e(0,12)$ & $1\text{--}8$ \\
4 & $e(8,12)$ & $1$ \\
5 & $e(1,12)$ & $5$ \\
6 & $v(5)$ & $1\text{--}10\text{--}11\text{--}9\text{--}12$ \\
7 & $v(10)$ & $1\text{--}13\text{--}3\text{--}6\text{--}11$ \\
8 & $v(3)$ & $6\text{--}12\text{--}9\text{--}13$ \\
9 & $v(6)$ & $11\text{--}7\text{--}2\text{--}12$ \\
10 & $v(2)$ & $7\text{--}13\text{--}9\text{--}12$ \\
11 & $v(7)$ & $11\text{--}8\text{--}1\text{--}13$ \\
12 & $v(1)$ & $8\text{--}0\text{--}13$ \\
13 & $v(0)$ & $8\text{--}4\text{--}13$ \\
14 & $v(4)$ & $8\text{--}11\text{--}9\text{--}13$ \\
15 & $v(8)$ & $11$ \\
16 & $v(11)$ & $9$ \\
17 & $v(12)$ & $9$ \\
18 & $v(9)$ & $13$ \\
\bottomrule
\end{tabular}

\end{table}

\begin{table}[htbp]
\centering\small
\caption{Reductions of the other two contractible seeds after one
admissible edge addition. A dash means that the singleton has already
been reached.}\label{tab:variant-reductions}
\begin{tabular}{@{}rllll@{}}
\toprule
& \multicolumn{2}{c}{$\Gb+\{0,3\}$} & \multicolumn{2}{c}{$\Gc+\{2,3\}$} \\
Step & Delete & Current neighbourhood & Delete & Current neighbourhood \\
\midrule
1 & $e(0,9)$ & $12\text{--}3\text{--}13$ & $e(2,7)$ & $10\text{--}3\text{--}13$ \\
2 & $e(0,13)$ & $3\text{--}7$ & $e(2,10)$ & $3\text{--}11$ \\
3 & $e(9,12)$ & $3$ & $e(2,13)$ & $3\text{--}6$ \\
4 & $e(3,9)$ & $13$ & $e(6,13)$ & $3$ \\
5 & $e(3,13)$ & $7$ & $e(3,13)$ & $7$ \\
6 & $e(7,13)$ & $2$ & $v(7)$ & $3\text{--}10\text{--}12\text{--}0\text{--}13$ \\
7 & $v(2)$ & $7\text{--}11\text{--}6\text{--}12\text{--}8\text{--}13$ & $v(10)$ & $3\text{--}11\text{--}5\text{--}1\text{--}12$ \\
8 & $v(6)$ & $11\text{--}4\text{--}10\text{--}3\text{--}12$ & $v(1)$ & $5\text{--}13\text{--}4\text{--}8\text{--}12$ \\
9 & $v(3)$ & $10\text{--}7\text{--}0\text{--}12$ & $v(4)$ & $8\text{--}11\text{--}9\text{--}0\text{--}13$ \\
10 & $v(0)$ & $7\text{--}5\text{--}12$ & $v(0)$ & $12\text{--}9\text{--}5\text{--}13$ \\
11 & $v(4)$ & $10\text{--}8\text{--}13\text{--}9\text{--}11$ & $v(8)$ & $11\text{--}2\text{--}6\text{--}12$ \\
12 & $v(7)$ & $10\text{--}5\text{--}11$ & $v(2)$ & $6\text{--}3\text{--}11$ \\
13 & $v(5)$ & $11\text{--}1\text{--}8\text{--}10,\ 8\text{--}12$ & $v(3)$ & $6\text{--}9\text{--}11$ \\
14 & $v(1)$ & $8\text{--}13\text{--}9\text{--}11$ & $v(6)$ & $9\text{--}12$ \\
15 & $v(10)$ & $8$ & $v(11)$ & $5\text{--}9$ \\
16 & $v(11)$ & $9$ & $v(12)$ & $9$ \\
17 & $v(9)$ & $13$ & $v(9)$ & $5$ \\
18 & $v(12)$ & $8$ & $v(5)$ & $13$ \\
19 & $v(8)$ & $13$ & --- & --- \\
\bottomrule
\end{tabular}

\end{table}
